\documentclass[11pt,a4paper]{article}

% ──────────────────────────────────────────────────────────────
% Packages
% ──────────────────────────────────────────────────────────────
\usepackage[utf8]{inputenc}
\usepackage[T1]{fontenc}
\usepackage{lmodern}
\usepackage[margin=2cm,top=2.5cm,bottom=2.5cm]{geometry}
\usepackage{graphicx}
\usepackage{xcolor}
\usepackage{hyperref}
\usepackage{amsmath,amssymb}
\usepackage{booktabs}
\usepackage{enumitem}
\usepackage{tabularx}
\usepackage{float}
\usepackage{fancyhdr}
\usepackage{titlesec}
\usepackage{microtype}
\usepackage{multicol}
\usepackage{caption}

% ──────────────────────────────────────────────────────────────
% Colors
% ──────────────────────────────────────────────────────────────
\definecolor{qcopper}{RGB}{184,135,100}
\definecolor{qdark}{RGB}{30,28,26}
\definecolor{qgold}{RGB}{207,181,141}
\definecolor{qsec}{RGB}{120,90,65}
\definecolor{linkblue}{RGB}{40,80,140}

% ──────────────────────────────────────────────────────────────
% Hyperlinks
% ──────────────────────────────────────────────────────────────
\hypersetup{
  colorlinks=true,
  linkcolor=qsec,
  citecolor=qsec,
  urlcolor=linkblue,
  pdftitle={QUILLON VAULT: A Post-Quantum Hardware Wallet with Physical Air-Gap Security},
  pdfauthor={Q-NarwhalKnight / Quillon},
}

% ──────────────────────────────────────────────────────────────
% Section styling
% ──────────────────────────────────────────────────────────────
\titleformat{\section}{\large\bfseries\color{qdark}}{\thesection.}{0.5em}{}[\vspace{-0.3em}\textcolor{qcopper}{\rule{\linewidth}{0.5pt}}]
\titleformat{\subsection}{\normalsize\bfseries\color{qsec}}{\thesubsection}{0.5em}{}

% ──────────────────────────────────────────────────────────────
% Header / Footer
% ──────────────────────────────────────────────────────────────
\pagestyle{fancy}
\fancyhf{}
\fancyhead[L]{\small\textcolor{qsec}{QUILLON VAULT}}
\fancyhead[R]{\small\textcolor{qsec}{Q-NarwhalKnight}}
\fancyfoot[C]{\small\thepage}
\renewcommand{\headrulewidth}{0.4pt}
\renewcommand{\headrule}{\hbox to\headwidth{\color{qcopper}\leaders\hrule height \headrulewidth\hfill}}

% ══════════════════════════════════════════════════════════════
\begin{document}
% ══════════════════════════════════════════════════════════════

% ── Title Page (single column) ──
\begin{center}
  {\Huge\bfseries\textcolor{qdark}{QUILLON VAULT}}\\[0.4em]
  {\Large\textcolor{qsec}{A Post-Quantum Hardware Wallet\\[0.15em]with Physical Air-Gap Security}}\\[0.6em]
  {\textcolor{qcopper}{\rule{0.5\textwidth}{0.8pt}}}\\[1em]
  {\large\textbf{Q-NarwhalKnight Project}}\\[0.3em]
  {\texttt{https://quillon.xyz}}\\[0.3em]
  {\small Version 1.0 --- February 2026}\\[1.2em]
  \textit{\large ``As thin as a credit card. As secure as a vault. As elegant as jewelry.''}
\end{center}

\vspace{0.8em}

% ── Full-width product image ──
\begin{figure}[H]
  \centering
  \includegraphics[width=\textwidth]{vault.png}
  \caption{QUILLON VAULT --- exploded component view (left), assembled device with
  OLED display and confirm/reject buttons (right), and USB-C slide mechanism
  in locked and exposed states (bottom insets).}
  \label{fig:vault}
\end{figure}

\vspace{0.5em}

% ── Abstract ──
\noindent\textcolor{qcopper}{\rule{\textwidth}{0.5pt}}
\vspace{0.3em}

\noindent\textbf{Abstract.}
We present \textbf{QUILLON VAULT}, a credit-card-form-factor hardware wallet
designed for the Q-NarwhalKnight post-quantum blockchain. At 85.6\,mm $\times$
54\,mm $\times$ 3.8\,mm and 28\,g, it is among the thinnest hardware signing
devices ever proposed. The defining innovation is a \emph{machined titanium
sliding frame} that physically disconnects the USB-C data lines when locked ---
creating a true hardware air gap that no software exploit can bridge. Internally,
a dedicated Quantum Random Number Generator (QRNG) chip provides
information-theoretically unpredictable entropy for key generation, while
NIST FIPS~204 Dilithium5 (ML-DSA, Level~5) signatures execute entirely
within a tamper-resistant secure element. Every transaction is dual-signed
with both Ed25519 and Dilithium5, providing security against both classical
and quantum adversaries. The USB channel is encrypted via Kyber-1024 (ML-KEM)
key encapsulation combined with AES-256-GCM. This paper details the
mechanical design, security architecture, cryptographic protocols, and
manufacturing considerations for the QUILLON VAULT.

\vspace{0.3em}
\noindent\textcolor{qcopper}{\rule{\textwidth}{0.5pt}}

\vspace{1em}

% ── Switch to two-column for the body ──
\begin{multicols}{2}

% ══════════════════════════════════════════════════════════════
\section{Introduction}
% ══════════════════════════════════════════════════════════════

The security of cryptocurrency hardware wallets rests on a simple principle:
\emph{private keys must never leave the device}. Yet existing wallets remain
vulnerable along two axes that current designs inadequately address.

\textbf{First}, the USB interface is a permanent attack surface. When a hardware
wallet is plugged in, the host computer can enumerate it, probe its descriptors,
and --- in the presence of firmware vulnerabilities --- extract secrets or inject
malicious payloads. Every minute a wallet remains connected is a minute of
exposure.

\textbf{Second}, all production hardware wallets today generate their keys
using deterministic pseudo-random number generators (PRNGs) seeded by
classical entropy sources. These are computationally secure under current
assumptions, but offer no information-theoretic guarantees. A sufficiently
powerful adversary --- including a future cryptanalytically relevant quantum
computer (CRQC) --- could potentially reconstruct the PRNG state.

The QUILLON VAULT addresses both problems with two hardware-level innovations:

\begin{enumerate}[leftmargin=1.5em,itemsep=2pt]
  \item A \textbf{physical air-gap mechanism}: a precision-machined titanium
    frame that slides along ceramic ball-bearing rails, actuating a hardware
    switch that electrically disconnects the USB-C D+/D$-$ data lines.
    When the frame is in the LOCKED position, the wallet is electrically
    unreachable --- not merely software-isolated, but circuit-broken.

  \item A \textbf{Quantum Random Number Generator} (QRNG) chip that derives
    entropy from quantum shot noise, providing keys whose randomness is
    guaranteed by the laws of quantum mechanics rather than computational
    hardness assumptions.
\end{enumerate}

These are complemented by a fully post-quantum cryptographic stack: Dilithium5
signatures, Kyber-1024 encrypted USB communication, and SHA-3-256 hashing.

% ══════════════════════════════════════════════════════════════
\section{Design Philosophy}
% ══════════════════════════════════════════════════════════════

The QUILLON VAULT was designed under three constraints:

\begin{description}[leftmargin=0.5em,itemsep=2pt]
  \item[Thinness.] The device must fit in a standard card slot. We target the
    ISO/IEC~7810 ID-1 footprint (85.6\,mm $\times$ 54\,mm) at a thickness of
    3.8\,mm --- thinner than two stacked credit cards.

  \item[Security by physics.] Every critical security property must have a
    physical enforcement mechanism, not merely a software one. Air gaps
    should be air gaps. Randomness should be quantum.

  \item[Elegance.] A security device that users find beautiful will be
    carried; one they find ugly will be left in a drawer. The titanium
    frame, sapphire glass, and precision slide mechanism are not cosmetic
    additions --- they are adoption enablers.
\end{description}

% ══════════════════════════════════════════════════════════════
\section{Mechanical Architecture}
% ══════════════════════════════════════════════════════════════

\subsection{Form Factor}

The device dimensions are summarized in Table~\ref{tab:dimensions}.
The thickness budget is tightly constrained:

\begin{center}
\small
\begin{tabular}{lr}
  \toprule
  \textbf{Layer} & \textbf{Thickness} \\
  \midrule
  Titanium top shell (Gr.~5) & 0.6\,mm \\
  Sapphire glass OLED window   & 0.3\,mm \\
  PCB + components (4-layer)   & 2.0\,mm \\
  Titanium bottom shell        & 0.6\,mm \\
  Magnetic slide rail           & 0.3\,mm \\
  \midrule
  \textbf{Total}               & \textbf{3.8\,mm} \\
  \bottomrule
\end{tabular}
\end{center}

\begin{table}[H]
  \centering
  \small
  \caption{Physical specifications.}
  \label{tab:dimensions}
  \begin{tabular}{ll}
    \toprule
    \textbf{Parameter} & \textbf{Value} \\
    \midrule
    Length  & 85.6\,mm \\
    Width   & 54\,mm \\
    Thickness & 3.8\,mm \\
    Weight (Ti) & 28\,g \\
    Weight (Al) & 22\,g \\
    Shell & Grade~5 titanium \\
    Display & Sapphire crystal \\
    PCB & 4-layer, 0.8\,mm \\
    \bottomrule
  \end{tabular}
\end{table}

\subsection{The Slide Mechanism}

The central mechanical innovation is the sliding titanium frame. An 8\,mm
linear translation exposes the USB-C port. The mechanism consists of:

\begin{itemize}[leftmargin=1.2em,itemsep=2pt]
  \item \textbf{Precision slide rails} machined into both top and bottom
    shells, with ceramic ball bearings (Si$_3$N$_4$) for zero-wobble
    linear travel.
  \item \textbf{N52 neodymium magnet detents} at both the LOCKED and
    UNLOCKED positions, providing a satisfying tactile ``click'' and
    holding the frame securely in either state.
  \item \textbf{A hardware disconnect switch}: a spring-loaded contact
    embedded in the slide rail that is physically depressed only when the
    frame reaches the UNLOCKED position. This switch completes the
    USB-C D+/D$-$ circuit. In the LOCKED position, the circuit is
    \emph{broken} --- an air gap enforced by mechanics, not software.
\end{itemize}

The electrical implication is critical: when the frame is closed, the
USB-C data lines are physically interrupted. No amount of software
exploitation on the host computer can establish communication with the
wallet. The power lines (VBUS/GND) are also disconnected, preventing
even side-channel power analysis via the USB connector.

\subsection{Exterior Finish}

All edges feature a 45$^\circ$ chamfer with mirror polish, contrasting
against the brushed face. Six material variants are planned:

\begin{enumerate}[leftmargin=1.2em,itemsep=1pt]
  \item \textbf{Obsidian} --- PVD black titanium, gold logo
  \item \textbf{Titanium} --- Natural Grade~5, silver logo
  \item \textbf{Rose} --- Rose gold PVD, white logo
  \item \textbf{Stealth} --- Matte black ceramic, no logo
  \item \textbf{Arctic} --- White zirconia ceramic, Ti frame
  \item \textbf{Carbon} --- Carbon fiber inlay, Ti frame
\end{enumerate}

% ══════════════════════════════════════════════════════════════
\section{Electronic Architecture}
% ══════════════════════════════════════════════════════════════

\subsection{Component Selection}

Table~\ref{tab:components} lists the key electronic components. Every part
was selected for the intersection of security properties, physical size,
and power characteristics compatible with a bus-powered, battery-less
design.

\begin{table}[H]
  \centering
  \small
  \caption{Key electronic components.}
  \label{tab:components}
  \begin{tabularx}{\columnwidth}{lX}
    \toprule
    \textbf{Part} & \textbf{Specification} \\
    \midrule
    SE & ATECC608B (2$\times$3\,mm) \\
    MCU & STM32L4S5, Cortex-M4, 120\,MHz, TrustZone \\
    Display & SSD1306 0.42'' OLED, 72$\times$40\,px \\
    USB & USB-C mid-mount (8.9$\times$2.5\,mm) \\
    QRNG & IDQ Quantis (3$\times$3\,mm) \\
    Flash & W25Q16JV, 2\,MB SPI NOR \\
    Haptics & LRA $\times$2 (4$\times$4$\times$1.5\,mm) \\
    Buttons & Capacitive touch (flush) \\
    \bottomrule
  \end{tabularx}
\end{table}

\subsection{Secure Element: ATECC608B}

The ATECC608B is a dedicated cryptographic coprocessor with
hardware-protected key storage (keys never leave the die),
side-channel resistance (DPA, SPA, fault injection),
a monotonic counter for replay attack prevention,
secure boot attestation, and a hardware TRNG
(supplemented by the external QRNG).

Private keys are generated \emph{inside} the secure element using entropy
from the QRNG chip and are used for signing without ever being exposed
to the MCU, the host computer, or firmware update processes.

\subsection{Quantum Random Number Generator}

The IDQ Quantis QRNG chip generates random bits from quantum vacuum
fluctuations (shot noise on a photodetector). Unlike PRNGs, the output
is \emph{information-theoretically random} --- its unpredictability is
guaranteed by quantum mechanics, not computational hardness. This
provides:

\begin{itemize}[leftmargin=1.2em,itemsep=2pt]
  \item Key generation entropy that cannot be predicted even by a
    quantum computer
  \item Nonce generation that eliminates the risk of
    nonce reuse via PRNG state compromise
  \item Seed material for the BIP-39 mnemonic that is physically
    unpredictable
\end{itemize}

\subsection{MCU: STM32L4S5 with TrustZone}

The STM32L4S5 provides ARM TrustZone hardware isolation, partitioning
the processor into a \textbf{Secure world} (cryptographic operations,
SE communication, signature construction) and a \textbf{Non-secure world}
(USB protocol, display rendering, user interaction).

A compromised non-secure world cannot access secure-world memory,
registers, or peripherals. Additional protections include stack canaries,
ASLR, and a hardware watchdog timer for fault injection resistance.

\subsection{Display and User Interaction}

The 0.42'' SSD1306 OLED display, protected by a 0.3\,mm sapphire
crystal window, shows transaction details (recipient, amount, token),
wallet addresses, firmware version, and PIN entry prompts.

Two flush-mount capacitive touch buttons (\textsc{Confirm} and
\textsc{Reject}) with LRA haptic feedback provide physical transaction
authorization. The haptic pattern differs between confirm
(single pulse) and reject (double pulse) for accessibility.

% ══════════════════════════════════════════════════════════════
\section{Security Architecture}
% ══════════════════════════════════════════════════════════════

The QUILLON VAULT implements defense-in-depth across five layers.

\subsection{Layer 1: Physical Security}

\begin{itemize}[leftmargin=1.2em,itemsep=2pt]
  \item \textbf{Titanium enclosure}: Grade~5 Ti-6Al-4V
    provides drill and cut resistance
  \item \textbf{Epoxy-potted PCB}: Optically opaque, thermally resistant
    encapsulation prevents physical probe access
  \item \textbf{Tamper mesh}: Fine copper mesh on the PCB
    detects intrusion and triggers key wipe
  \item \textbf{Holographic seal}: Tamper-evident seal
    over the shell seam
\end{itemize}

\subsection{Layer 2: Secure Element Isolation}

The ATECC608B enforces that key material never crosses the die boundary,
all signing operations execute within the SE, side-channel countermeasures
operate at the silicon level, and a monotonic counter prevents replay.

\subsection{Layer 3: Firmware Integrity}

Secure boot chain requires firmware signed by Quillon's release key
(stored in SE). TrustZone memory isolation prevents cross-world leakage.
Stack canaries and ASLR harden against memory corruption. A hardware
watchdog detects voltage glitching and clock manipulation.

\subsection{Layer 4: Post-Quantum Cryptography}
\label{sec:crypto}

\begin{description}[leftmargin=0.5em,itemsep=3pt]
  \item[Transaction signatures:] \textbf{Dual-signature mode.} Every
    transaction is signed with both Ed25519 (classical, 64-byte)
    and Dilithium5 / ML-DSA (post-quantum, NIST Level~5,
    4,627-byte). If a CRQC breaks Ed25519, Dilithium5
    remains valid. If lattice cryptanalysis advances, Ed25519 remains
    secure. Both must verify.

  \item[Key generation:] QRNG-derived entropy, processed through
    SHA-3-256, provides seed material for both key pairs.
    QRNG output is XORed with the SE's internal TRNG
    as a defense-in-depth entropy source.

  \item[USB encryption:] Host-to-wallet communication uses
    \textbf{Kyber-1024} (ML-KEM, FIPS~203) for
    key encapsulation, establishing a shared secret for
    \textbf{AES-256-GCM} authenticated encryption. Even a
    compromised USB stack cannot extract signing keys.

  \item[Hashing:] \textbf{SHA-3-256} (FIPS~202) for all internal
    hashing. SHA-3-256 provides $\sim$128-bit security against
    Grover's quantum search.
\end{description}

\subsection{Layer 5: User Verification}

The OLED displays \emph{exact} transaction details --- what the
user sees is what the SE signs. A physical button press is required;
remote signing is impossible. PIN policy: 3 wrong attempts triggers
a 24-hour lockout; 10 wrong attempts triggers permanent key wipe.
Recovery via BIP-39 mnemonic on titanium seed card.

% ══════════════════════════════════════════════════════════════
\section{The Physical Air Gap}
% ══════════════════════════════════════════════════════════════

The slide mechanism's security properties deserve emphasis.
In the \textbf{LOCKED} state:

\begin{equation}
  \text{USB D+/D}^{-} \;\longrightarrow\; \underbrace{[\;\text{GAP}\;]}_{\text{switch open}} \;\longrightarrow\; \text{MCU}
\end{equation}

The spring-loaded contact in the slide rail is mechanically retracted,
creating a physical break in the D+ and D$-$ traces. No electrical
signal can traverse this gap. The VBUS power line is similarly
disconnected, preventing USB enumeration, firmware exploitation,
side-channel power analysis, and EM emanation analysis via the USB
cable as antenna.

In the \textbf{UNLOCKED} state (frame slid 8\,mm), the spring contact
is depressed, completing the circuit:

\begin{equation}
  \text{USB D+/D}^{-} \;\xrightarrow{\;\text{contact}\;}\; \text{MCU}
\end{equation}

This is, to our knowledge, the first hardware wallet where the
USB data connection has a \emph{user-controlled physical disconnect}
as an integral part of the device form factor --- not an afterthought
(USB cap) or an accessory (Faraday bag), but a machined precision
mechanism built into the device's primary interaction paradigm.

% ══════════════════════════════════════════════════════════════
\section{Cryptographic Protocol}
% ══════════════════════════════════════════════════════════════

\subsection{Key Generation}

\begin{enumerate}[leftmargin=1.2em,itemsep=2pt]
  \item QRNG chip generates 512 bits of quantum-random entropy
  \item SE internal TRNG generates 512 bits of classical entropy
  \item Both sources are XORed, then processed through SHA-3-512
  \item Output seeds both Ed25519 and Dilithium5 key pairs
  \item BIP-39 mnemonic derived from the combined 512-bit seed
  \item Key pairs generated \emph{inside the SE}, never exported
\end{enumerate}

\subsection{Transaction Signing}

\begin{enumerate}[leftmargin=1.2em,itemsep=2pt]
  \item Host sends transaction over Kyber-1024 encrypted channel
  \item MCU renders transaction details on OLED display
  \item User taps \textsc{Confirm} or \textsc{Reject}
  \item If confirmed: MCU passes transaction hash to SE
  \item SE signs with Ed25519, then with Dilithium5
  \item Both signatures returned over encrypted channel
  \item SE increments monotonic counter (prevents replay)
\end{enumerate}

\subsection{Firmware Updates}

New firmware images must be signed with Quillon's Dilithium5 release key.
The SE verifies the signature before allowing flash writes. The OLED
displays the version change for user confirmation, and a physical button
press is required to authorize. The previous firmware is retained in a
backup partition for rollback, and the update process is atomic
(power-loss-safe).

% ══════════════════════════════════════════════════════════════
\section{Threat Model}
% ══════════════════════════════════════════════════════════════

The following threats are addressed:

\smallskip
\noindent\textbf{Malware on host.} USB exploitation, firmware injection.
\textit{Mitigated by:} Physical air gap; Kyber-1024 encrypted channel;
signed firmware updates.

\smallskip
\noindent\textbf{Physical theft.} Device stolen, keys extracted.
\textit{Mitigated by:} PIN lockout; titanium shell; epoxy-potted PCB;
tamper mesh triggers key wipe.

\smallskip
\noindent\textbf{Side-channel.} Power analysis, EM emanation, timing.
\textit{Mitigated by:} SE with DPA/SPA resistance; air gap disconnects
power; constant-time crypto.

\smallskip
\noindent\textbf{Supply chain.} Tampered device before delivery.
\textit{Mitigated by:} Holographic seal; secure boot attestation;
device certificate chain.

\smallskip
\noindent\textbf{Quantum adversary.} Shor's algorithm on ECDSA/EdDSA.
\textit{Mitigated by:} Dual-sign Dilithium5 + Ed25519; QRNG keys;
Kyber-1024 USB encryption.

\smallskip
\noindent\textbf{PRNG weakness.} Predictable key generation.
\textit{Mitigated by:} QRNG chip (quantum shot noise); XOR with SE TRNG;
information-theoretic randomness.

\smallskip
\noindent\textbf{Evil maid.} Device tampered while unattended.
\textit{Mitigated by:} Tamper mesh; holographic seal; secure boot detects
firmware modification.

% ══════════════════════════════════════════════════════════════
\section{Comparison with Existing Wallets}
% ══════════════════════════════════════════════════════════════

\begin{table}[H]
  \centering
  \small
  \caption{Feature comparison.}
  \label{tab:comparison}
  \begin{tabularx}{\columnwidth}{lX}
    \toprule
    \textbf{Feature} & \textbf{Status vs.\ Competitors} \\
    \midrule
    Thickness & \textbf{3.8\,mm} (vs.\ 8--15\,mm) \\
    Physical air gap & \textbf{USB slide disconnect} (unique) \\
    PQ signatures & \textbf{Dilithium5 Level~5} (none offer PQ) \\
    QRNG & \textbf{IDQ Quantis} (none offer QRNG) \\
    Dual-sign & \textbf{Ed25519 + Dilithium5} (unique) \\
    Encrypted USB & \textbf{Kyber-1024 + AES-256} (unique) \\
    Secure Element & ATECC608B (comparable) \\
    Shell material & \textbf{Grade~5 Titanium} (vs.\ plastic) \\
    \bottomrule
  \end{tabularx}
\end{table}

No existing production hardware wallet offers post-quantum signatures,
quantum random number generation, or a physical USB disconnect mechanism.
The QUILLON VAULT is, to our knowledge, the first design to combine all
three in a credit-card form factor.

% ══════════════════════════════════════════════════════════════
\section{User Experience Flow}
% ══════════════════════════════════════════════════════════════

The interaction model prioritizes simplicity:

\begin{enumerate}[leftmargin=1.2em,itemsep=3pt]
  \item \textbf{Slide frame open} --- magnetic click at UNLOCKED
  \item \textbf{Insert USB-C} --- OLED wakes, shows wallet address
  \item \textbf{Authorize transaction} --- OLED shows amount,
    recipient, token; user taps \textsc{Confirm} or \textsc{Reject}
  \item \textbf{Unplug USB-C} --- OLED enters sleep
  \item \textbf{Slide frame closed} --- magnetic click at LOCKED;
    USB-C electrically disconnected and physically concealed
\end{enumerate}

The entire signing interaction takes approximately 5 seconds.
The device requires no battery --- it is entirely bus-powered via USB.

% ══════════════════════════════════════════════════════════════
\section{Manufacturing and Cost}
% ══════════════════════════════════════════════════════════════

\begin{table}[H]
  \centering
  \small
  \caption{Bill of materials (10K volume).}
  \label{tab:bom}
  \begin{tabular}{lr}
    \toprule
    \textbf{Component} & \textbf{Cost} \\
    \midrule
    Titanium shells (CNC + PVD)   & \$18.00 \\
    PCB + assembly (4-layer)       & \$8.00 \\
    ATECC608B Secure Element       & \$1.20 \\
    STM32L4S5 MCU                  & \$6.00 \\
    0.42'' OLED + sapphire         & \$4.00 \\
    USB-C mid-mount connector      & \$0.80 \\
    IDQ Quantis QRNG               & \$12.00 \\
    Slide mechanism                & \$5.00 \\
    LRA haptic motors ($\times$2)   & \$1.50 \\
    Flash + passives + assembly    & \$3.00 \\
    \midrule
    \textbf{Total BOM}             & \textbf{\$59.50} \\
    \midrule
    Retail (Titanium)              & \$149.00 \\
    Retail (Ceramic)               & \$199.00 \\
    \bottomrule
  \end{tabular}
\end{table}

The titanium shells are 5-axis CNC machined with PVD coating.
The slide mechanism uses injection-molded rail channels
with press-fit ceramic ball bearings and N52 neodymium magnets epoxied
into machined pockets.

% ══════════════════════════════════════════════════════════════
\section{Packaging}
% ══════════════════════════════════════════════════════════════

The QUILLON VAULT ships in a magnetic-close presentation box with
soft-touch matte finish and foil-stamped Q logo:

\begin{itemize}[leftmargin=1.2em,itemsep=2pt]
  \item Device in molded microfiber tray
  \item Braided USB-C cable (30\,cm)
  \item \textbf{Titanium recovery seed card} --- fireproof, waterproof,
    for metal-stamping the 24-word BIP-39 mnemonic
  \item Quick-start guide with QR code
\end{itemize}

The titanium seed card ensures the backup survives conditions that would
destroy paper or plastic.

% ══════════════════════════════════════════════════════════════
\section{Integration with Q-NarwhalKnight}
% ══════════════════════════════════════════════════════════════

The QUILLON VAULT is the native hardware signing device
for the Q-NarwhalKnight post-quantum blockchain~\cite{qnk2026}. The
integration is natural because Q-NarwhalKnight already
requires dual Ed25519 + Dilithium5 signatures, SHA-3-256 transaction
hashing, and Kyber-1024 for P2P encryption.

Supported operations include native QUG transfers, QUGUSD stablecoin
transactions, custom token (QRC-20) transfers, DEX swap authorization,
smart contract interaction, and RWA token management.

% ══════════════════════════════════════════════════════════════
\section{Future Work}
% ══════════════════════════════════════════════════════════════

\begin{itemize}[leftmargin=1.2em,itemsep=2pt]
  \item \textbf{Bluetooth Low Energy}: Wireless signing with
    Kyber-1024 encrypted channel
  \item \textbf{NFC tap-to-sign}: Contactless transaction
    authorization at point-of-sale
  \item \textbf{Multi-signature}: $m$-of-$n$ threshold
    signing across multiple devices
  \item \textbf{Biometric auth}: Fingerprint sensor
    with SE-backed template storage
  \item \textbf{SQIsign migration}: When isogeny-based signatures
    mature, the crypto-agile firmware allows switching
    from Dilithium5 to SQIsign for 204-byte PQ signatures
\end{itemize}

% ══════════════════════════════════════════════════════════════
\section{Conclusion}
% ══════════════════════════════════════════════════════════════

The QUILLON VAULT demonstrates that post-quantum hardware security,
physical air-gap isolation, and elegant industrial design are not
mutually exclusive. By building the security architecture
around physical enforcement mechanisms --- a machined titanium slide
that breaks USB circuits, a QRNG chip that derives entropy from
quantum mechanics, and a secure element that performs Dilithium5
signing without ever exposing key material --- we achieve security
properties independent of software correctness assumptions.

The device fits in a card slot at 3.8\,mm, weighs 28\,g,
and requires no battery. It is the first hardware wallet designed from
the ground up for the post-quantum era, and the first where the USB
connection has a user-controlled physical disconnect built into the
device's primary form factor.

\vspace{1em}
\noindent\textcolor{qcopper}{\rule{\columnwidth}{0.8pt}}

\vspace{0.5em}
\noindent\small\textit{The QUILLON VAULT is designed for Q-NarwhalKnight.
Visit} \texttt{https://quillon.xyz}.

% ══════════════════════════════════════════════════════════════
% References
% ══════════════════════════════════════════════════════════════

\begin{thebibliography}{99}

\bibitem{qnk2026}
  Q-NarwhalKnight Project,
  ``Q-NarwhalKnight: A Post-Quantum DAG Consensus System,''
  Technical whitepaper, 2026.

\bibitem{fips204}
  NIST,
  ``FIPS 204: Module-Lattice-Based Digital Signature Standard (ML-DSA),''
  2024.

\bibitem{fips203}
  NIST,
  ``FIPS 203: Module-Lattice-Based Key-Encapsulation Mechanism Standard (ML-KEM),''
  2024.

\bibitem{fips202}
  NIST,
  ``FIPS 202: SHA-3 Standard,''
  2015.

\bibitem{dilithium}
  L. Ducas et al.,
  ``CRYSTALS-Dilithium: A Lattice-Based Digital Signature Scheme,''
  TCHES, 2018.

\bibitem{kyber}
  R. Avanzi et al.,
  ``CRYSTALS-Kyber: A CCA-Secure Module-Lattice-Based KEM,''
  IEEE Euro S\&P, 2018.

\bibitem{ed25519}
  D.J. Bernstein et al.,
  ``High-Speed High-Security Signatures,''
  J. Cryptographic Engineering, 2012.

\bibitem{atecc608b}
  Microchip Technology,
  ``ATECC608B CryptoAuthentication Device Datasheet,''
  2020.

\bibitem{qrng}
  ID Quantique,
  ``Quantis QRNG Chip: True Random Number Generation,''
  2023.

\bibitem{bip39}
  M. Palatinus and P. Rusnak,
  ``BIP-39: Mnemonic Code for Generating Deterministic Keys,''
  2013.

\end{thebibliography}

\end{multicols}

\end{document}
